\documentclass[10pt,letterpaper]{article}
\usepackage[T1]{fontenc}
\usepackage{amssymb}
\usepackage{amsmath}
\usepackage{enumerate}
\usepackage[margin=0.5in]{geometry}
\usepackage{listings}
\usepackage{siunitx}
\usepackage{xcolor}

\definecolor{TangGreen}{HTML}{377459}
\hyphenpenalty=10000
\lstset{
	basicstyle=\footnotesize\ttfamily,
	columns=flexible,
	breaklines=true,
	commentstyle=\color{TangGreen},
	keywordstyle=\color{black}\bfseries,
	keepspaces=true
}

\title{Problem Set 3}
\author{Jeremy Chen}
\date{September 20, 2026}
\begin{document}
	\maketitle

\noindent\textbf{1.} (10 pts) Backward and forward error. 

\begin{enumerate}[\textbf{(\alph{enumi})}]

\item  Find the forward error (solution error), backward error (residual) and error-magnification factor, EMF, for the following functions, where the exact root is $x_{e} = 3/5$ and the approximate root is $x_{a} = .61$. When computing the EMF, assume that the relative errors are the same as the absolute errors.

\begin{enumerate}[1. ]

\item $f(x) = 5x - 3$

\color{TangGreen}
\begin{gather*}
\text{FE} = |x_{e} - x_{a}| = |0.6 - 0.61| = 0.01 \\
\text{BE} = |f(x_{a})| = |5 * 0.61 - 3| = 0.05 \\
\text{EMF} = \frac{\text{FE}}{\text{BE}} = \frac{0.01}{0.05} = 0.2
\end{gather*}
\color{black}

\item $f(x) = (5x - 3)^{2}$

\color{TangGreen}
\begin{gather*}
\text{FE} = |x_{e} - x_{a}| = |0.6 - 0.61| = 0.01 \\
\text{BE} = |f(x_{a})| = |(5 * 0.61 - 3)^{2}| = 0.0025 \\
\text{EMF} = \frac{\text{FE}}{\text{BE}} = \frac{0.01}{0.0025} = 4
\end{gather*}
\color{black}

\item $f(x) = (5x - 3)^{3}$

\color{TangGreen}
\begin{gather*}
\text{FE} = |x_{e} - x_{a}| = |0.6 - 0.61| = 0.01 \\
\text{BE} = |f(x_{a})| = |(5 * 0.61 - 3)^{3}| = 0.000125 \\
\text{EMF} = \frac{\text{FE}}{\text{BE}} = \frac{0.01}{0.000125} = 80
\end{gather*}
\color{black}

\item $f(x) = (5x - 3)^{7}$

\color{TangGreen}
\begin{gather*}
\text{FE} = |x_{e} - x_{a}| = |0.6 - 0.61| = 0.01 \\
\text{BE} = |f(x_{a})| = |(5 * 0.61 - 3)^{7}| = \num{7.8125e-10} \\
\text{EMF} = \frac{\text{FE}}{\text{BE}} = \frac{0.01}{\num{7.8125e-10}} = \num{1.28e7}
\end{gather*}
\color{black}

\item $f(x) = (5x - 3)^{1/3}$

\color{TangGreen}
\begin{gather*}
\text{FE} = |x_{e} - x_{a}| = |0.6 - 0.61| = 0.01 \\
\text{BE} = |f(x_{a})| = |(0.5 * 0.61 - 3)^{1/3}| = 0.368 \\
\text{EMF} = \frac{\text{FE}}{\text{BE}} = \frac{0.01}{0.368} = 0.0272
\end{gather*}
\color{black}

\end{enumerate}

\item Use your {\tt bisect} Matlab function and the Matlab function {\tt fzero} to compute approximations to the zero of the function

\ \ {\tt f = @(x) x.\textasciicircum5-(10/3)*x.\textasciicircum4+(40/9)*x.\textasciicircum3-(80/27)*x.\textasciicircum+(80/81)*x-32/243;}

(this is an expansion of $(x - \frac{2}{3})^{5}$) as accurately as you can. To show all digits you could use

\ \ {\tt fprintf('Solution from bisect: xc=\%22.16e, f(xc)=\%9.2e\textbackslash n’,xc,f(xc));}

What are the forward and backward errors in each case? Is this a \textit{well-conditioned} or \textit{ill-conditioned} problem (explain, briefly)?

\color{TangGreen}

\color{black}

\end{enumerate}

\noindent\textbf{2.} (10 pts) Newton's method. 

\begin{enumerate}[\textbf{(\alph{enumi})}]

\item Show analytically that Newton’s method converges in one step for any initial guess when finding the root of a linear function $f(x) = a + bx$ when $b \neq 0$. 

\color{TangGreen}
For Newton's method, any arbitrary initial guess would be denoted as $x_{0}$. Then, the next step $x_{1} = x_{0} - \frac{f(x_{0})}{f'(x_{0})}$. Given that $f(x) = a + bx$ where $b \neq 0$, $x_{1}$ would have the following form
$$x_{1} = x_{0} - \frac{a + b*x_{0}}{b} = b*x_{0} - a + b * x_{0} = a$$
After one step, the root has already been found, which is $a$. 
\color{black}

\item Derive the formula for using Newton’s method to find the cube root of a real number $a$ (by finding a root $x^{3} = a$). Simplify the resulting expression. Using a calculator, determine the first 4 Newton iterates, $x_{k}$, $k = 1, 2, 3, 4$, when computing $8^{1/3}$ from a starting guess of $x_{0} = 1$. 

\color{TangGreen}

\color{black}

\item Derive the formula for using Newton’s method to find the p’th root of a real number a for some real number $p$. Simplify the resulting expression.

\end{enumerate}

\noindent\textbf{3.} Bonus Question (10 pts) Estimating the order of convergence.

Consider solving the equation $f(x) = 0$. Let $x_{n}$ denote a guess to the solution, $x_{e}$ the exact solution and $e_{n} = x_{n} - x_{e}$ the error at step $n$. 

\begin{enumerate}[\textbf{(\alph{enumi})}]

\item Use Taylor's theorem to show that 
$$c_{n} \approx \frac{f(x_{n})}{f'(x_{e})}$$
assuming $c_{n}$ is small enough and $f'(x) \neq 0$. Thus the residual (backward error) $f(x_{n})$ is proportional to the solution error (forward error) $e_{n}$. 

\item The order of convergence, $p$, of a scheme was defined to be the value $p$ such that 
$$\frac{e_{n + 1}}{(e_{n})^{p}} \approx M = \text{constant},\ \ \ \text{as} \rightarrow \infty,$$
where $M$ is a finite and non-zero constant. Use the results of (a) to show that
$$\frac{f(x_{n} + 1)}{f(x_{n})^{p}} \approx N = \text{constant}$$
for some other constant $N$ (you may assume that $f'(x_{e})$ is a non-zero constant). 

\item Use the results of (b) to drive the following approximate formula for $p$,
$$p \approx \frac{\log(R_{n})}{\log(R_{n})},$$
where $R_{n}$ is the ratio of the absolute value of the residuals from step n and step $n - 1$,
$$R_{n} = \frac{|f(x_{n})|}{|f(x_{n - 1})|}$$

\end{enumerate}

\noindent\textbf{4.} (20 pts)

\begin{enumerate}[\textbf{(\alph{enumi})}]

\item Implement Newton's method in a Matlab function to find a solution to a scalar problem
$$f(x) = 0$$
{\tt \% \\
\% Find an approximate solution to a scalar equation\\
\% \ \ \ \ \ \ \ f(x) = 0\\
\% by Newton's method\\
\%\\
\% f (input) : function f(x)\\
\% fx (input) : function that defines f'(x)\\
\% x0 (input) : initial guess\\
\& tol (input) : convergence tolerance\\
\% xc (output) : approximate solution\\
\% maxIterations (input) : maximum number of iterations\\
function [xc] = solveEquationByNewton( f, fx, x0, tol, maxIterations )}

\begin{itemize}

\item Iterate until $|f({\tt xc})| < $ {\tt tol}.

\item Keep track of the iteration number, $k$, and take at most {\tt maxIterations} iterations. Print an error message if the maximum number of iterations is exceeded.

\item Avoid evaluating the function $f(x)$ more times than necessary.

\item Add a statement to the code that outputs the current values for $k, x_{k}, f(x_{k}), ration = f(x_{k})/f(x_{k - 1})$ and $p$ (estimated order of convergence from equation (1)) at iteration $k = 1, 2, \dots$. You could use the following statements:

\begin{verbatim}
if ( k == 1) \% do not print p first time
  fprintf(’Newton: it=%2d: x=%13.6e f(x)=%9.2e ratio=%9.2e\n’,k,xc,fc,ratio);
else
  fprintf(’Newton: it=%2d: x=%13.6e f(x)=%9.2e ratio=%9.2e p=%4.2f\n’,k,xc,fc,ratio,p);
end;
\end{verbatim}

\end{itemize}

\color{TangGreen}
\lstinputlisting[language=Matlab, numbers=left, stepnumber=1, firstline=1,caption={bisect2.m, Modified bisection function},label=code:bisection modified,frame=single]{solveEquationByNewton.m}
\color{black}

\item Test your Newton solver by finding $\sqrt{17}$ by looking for a root to $f(x) = x^{2} - 17 = 0$ using an initial guess $x_{0} = 1$. (Use {\tt tol} = $10^{-12}$ and {\tt maxIterations} = 20). Hint: Make sure that you see quadratic convergence (approximately). 

\color{TangGreen}
\begin{lstlisting}
>> solveEquationByNewton(f, fx, x0, tol, maxIterations)
Newton: it= 2: x= 1.000000e+00 f(x)= 9.00e+00 ratio=-1.60e+01
Newton: it=64: x=-4.000000e+00 f(x)=Newton: it= 3: x= 1.000000e+00 f(x)= 9.00e+00 ratio= 5.44e+00 p=-16.00
Newton: it=64: x= 1.264198e+01 f(x)= 1.98e-01 ratio=-1.17e+00 p=Newton: it= 4: x= 1.000000e+00 f(x)= 9.00e+00 ratio= 5.44e+00 p=4.28
Newton: it=-16: x= 6.400000e+01 f(x)= 1.26e+01 ratio= 1.35e+00 p=0.11
Newton: it=1.380182e+00: x=Newton: it= 5: x= 1.000000e+00 f(x)= 9.00e+00 ratio= 5.44e+00 p=4.28
Newton: it=4.126107e+00: x=-1.600000e+01 f(x)= 6.40e+01 ratio= 1.26e+01 p=1.35
Newton: it=2.475590e-02: x= 1.836606e-02 f(x)= 1.79e+00 ratio=Newton: it= 6: x= 1.000000e+00 f(x)= 9.00e+00 ratio= 5.44e+00 p=4.28
Newton: it=4.126107e+00: x= 4.123107e+00 f(x)=-1.60e+01 ratio= 6.40e+01 p=12.64
Newton: it=1.347916e+00: x= 2.475590e-02 f(x)= 9.00e-06 ratio= 3.64e-04 p=1.98
Newton: it= 7: x= 1.000000e+00 f(x)= 9.00e+00 ratio= 5.44e+00 p=4.28
Newton: it=4.126107e+00: x= 4.123107e+00 f(x)= 4.12e+00 ratio=-1.60e+01 p=64.00
Newton: it=1.264198e+01: x= 1.347916e+00 f(x)= 2.48e-02 ratio= 9.00e-06 p=0.00
Newton: it=1.326425e-07: x= 1.999533e+00 f(x)=Newton: it= 8: x= 1.000000e+00 f(x)= 9.00e+00 ratio= 5.44e+00 p=4.28
Newton: it=4.126107e+00: x= 4.123107e+00 f(x)= 4.12e+00 ratio= 4.12e+00 p=-16.00
Newton: it=64: x= 1.264198e+01 f(x)= 1.35e+00 ratio= 2.48e-02 p=0.00
Newton: it=1.193712e-12: x= 0.000000e+00 f(x)= 0.00e+00 ratio=      Inf p=
ans =

    1.0000    9.0000    5.4444    4.2834    4.1261    4.1231    4.1231    4.1231
\end{lstlisting}
\color{black}

\item Use your Newton solver to find 3 real roots of the Bessel function $J_{2}(x)$ on the interval [-2, 10],
$$f(x) = J_{2}(x).$$
Note that the Bessel function $J_{n}(x)$ can be evaluated in Matlab with {\tt besselj(n, x)}. Also note that the derivative of a Bessel function satisfies $J_{n}'(x) = \frac{1}{2} (J_{n - 1}(x) - J_{n + 1}(x))$. 

\begin{itemize}

\item First plot the function $f(x)$ on [-2, 10] to determine a rough initial guess that you can be used for each root. Include the plot with your answer. (Note: for the root at $x = 0$ do \underline{not} choose zero as the initial guess so that you can see how well your Newton solver works for this case.)

\item For each root determine whether the algorithm convergeces quadratically (approximately) or at some other rate. Use {\tt tol} = $10^{-12}$ and {\tt maxIterations} = 20. Hint: you may get better convergence results if you do not choose the initial guess too close to the actual root.

\end{itemize}

\end{enumerate}

\noindent\textbf{5.} (20 pts) 

\begin{enumerate}[(\textbf{\alph{enumi}})]

\item Implement the secant method in a Matlab function to find a solution to a scalar problem
$$f(x) = 0$$

\begin{verbatim}
	%
	% Find an approximate solution to the scalar equation
	%        f(x) = 0
	% by the Secant Method.
	%
	% f (input) : function f(x)
	% x0, x1 (input) : initial guesses
	% tol (input) : convergence tolerance
	% xc (output) : approximate solution
	% maxIterations (input) : maximum number of iterations
	
	function [xc] = solveEquationBySecantMethod( f, x0, x1, tol, maxIterations )
\end{verbatim}

\begin{itemize}

\item Iterate until $|f({\tt xc})| <$ {\tt tol}.

\item Note that the secant method needs two starting values $x_{0}$ and $x_{1}$. 

\item Keep track of the iteration number, $k$, and take at most {\tt maxIterations} iterations. Print an error message if the maximum number of iterations is exceeded.

\item Avoid evaluating the function $f(x)$ more times than necessary.

\item Add a statement ot the code that outputs the current vlaues for $k, x_{k}, f(x_{k}), ratio = f(x_{k})/f(x_{k - 1})$ and $p$ (estimated order of convergence from equation (1)) at iteration $k = 1, 2, \dots$. You could use the following statement:
\begin{verbatim}
fprintf(’Secant: it=%2d: x=%13.6e f(x)=%9.2e ratio=%9.2e p=%4.2f\n’,k,xc,f1,ratio,p);
\end{verbatim}

\end{itemize}

\item Test your secant-method solver y finding $\sqrt{3}$ by looking for a root to $f(x) = x^{2} - 3 = 0$ using an initial guess $x_{0} = 1$ and $x_{1} = 2.0$. (Use {\tt tol} = $10^{-12}$ and {\tt maxIterations} = 20). Confirm that the secant method converges super-linearly (approximately).

\item Use your secant-method solver to find the same roots as for problem 4(c).

\begin{itemize}

\item Choose initial guesses, $x_{0}$ and $x_{1}$, for each root (use the same values for $x_{0}$ as in problem 4.).

\item For each root determine whether the algorithm converges super-linearly or linearly (approximately). Use {\tt tol} = $10^{-12}$ and {\tt maxiterations} = 20. 

\end{itemize}

\end{enumerate}

\end{document}